\chapter{Nail Discoloration}

\section*{Definition}
Nail discoloration (chromonychia) refers to abnormal color of the nail plate or nail bed, which may reflect exogenous staining, benign conditions, or systemic disease. Patterns include leukonychia (white), melanonychia (brown/black), erythronychia (red), chloronychia (green), and xanthonychia (yellow).

\section*{What Happens in Your Body}
Nail plate color derives from the transparency of the nail, pigmentation of the underlying nail bed, and the nail plate substance itself. Leukonychia results from abnormal keratinization with retained air pockets within the nail plate (true) or from nail bed pallor (apparent). Melanonychia involves increased melanin deposition from melanocyte activation (normal number) or melanocytic hyperplasia (increased number, including nevus and melanoma). Yellow nails reflect reduced growth rate and thickened nail plates from lymphatic stasis or extrinsic staining.

\section*{Causes}
\begin{itemize}
  \item \textbf{Leukonychia (white):} Trauma, zinc deficiency, cirrhosis, chemotherapy, arsenic poisoning (Mees lines), fungal onychomycosis, chronic renal failure
  \item \textbf{Melanonychia (brown/black):} Racial pigmentation, trauma, fungal infection (Trichophyton rubrum), drugs (minocycline, chemotherapeutics, antimalarials), systemic lupus, Addison's disease, acral lentiginous melanoma
  \item \textbf{Chloronychia (green):} \textit{Pseudomonas aeruginosa} colonization (pyocyanin pigment); common in chronic paronychia and onycholysis
  \item \textbf{Xanthonychia (yellow):} Onychomycosis, yellow nail syndrome (lymphedema, pleural effusion), smoking, topical agents (tetanus toxoid, dithranol)
  \item \textbf{Erythronychia (red):} Trauma, glomus tumor, Darier disease, linear/lichen planus, splinter hemorrhages (endocarditis, trauma)
\end{itemize}

\section*{When to See a Doctor}
See your doctor if:
\begin{itemize}
  \item Longitudinal pigmented band with irregular borders, variegated color, or Hutchinson sign (pigment extension onto proximal nail fold) --- suspect subungual melanoma
  \item Diffuse nail discoloration with systemic symptoms: weight loss, dyspnea, edema, fatigue
  \item Green-black discoloration with associated paronychia or pain
  \item New or changing pigmentation in a single digit
\end{itemize}

\section*{Related Symptoms}
\begin{itemize}
  \item Onychomycosis, melanonychia striata, splinter hemorrhages, Muehrcke lines, Terry nails, Lindsay nails
\end{itemize}
\textbf{Important:} The ABCDEF mnemonic for subungual melanoma: A (age >60, African/Asian/Hispanic), B (brown/black band), C (change), D (digit involved >3 mm), E (extension), F (family history).

\section*{Self-Care / Home Management}
\begin{itemize}
  \item Keep fingernails clean and dry; avoid trauma and nail biting
  \item Trim nails straight across and avoid tight shoes to prevent trauma-induced discoloration
  \item Discontinue smoking, which stains nails and contributes to yellow discoloration
\end{itemize}

\section*{How Your Doctor Will Evaluate You}
\begin{itemize}
  \item Detailed history: onset, progression, trauma, medications, occupation, systemic symptoms
  \item Dermoscopy of the nail plate and nail bed for pigment pattern analysis (lines, dots, blotches)
  \item KOH preparation and fungal culture for suspected onychomycosis
  \item Nail plate biopsy or matrix biopsy for persistent or suspicious pigmented bands
\end{itemize}

\section*{Treatment Options}
\begin{itemize}
  \item Fungal: Oral terbinafine 250 mg daily for 6--12 weeks (fingernails) or 12--24 weeks (toenails)
  \item Pseudomonas-associated chloronychia: Topical gentamicin or ciprofloxacin drops, dilute bleach soaks
  \item Melanocytic: Surgical biopsy/excision of suspicious pigmented lesions; melanoma requires wide excision and sentinel node evaluation
  \item Systemic: Treat underlying disease (Addison's, renal failure, liver disease); discontinue offending medications
  \item Yellow nail syndrome: Vitamin E 1000 IU/day, intralesional steroids, topical antifungals
\end{itemize}

\section*{References}
\begin{thebibliography}{99}
\bibitem{nail_discoloration:ref1} Scher RK, Rich P, Pariser D, et al. "Nail disorders in dermatology." \textit{Journal of the American Academy of Dermatology}, 2014;71(2):373--383.
\bibitem{nail_discoloration:ref2} Haneke E. "Nail disorders in systemic diseases." \textit{Dermatologic Clinics}, 2006;24(3):341--347.
\bibitem{nail_discoloration:ref3} Tosti A, Piraccini BM. "Treatment of common nail disorders." \textit{Dermatologic Clinics}, 2000;18(3):407--416.
\bibitem{nail_discoloration:ref4} Baran R, Dawber RPR, Haneke E, et al. \textit{Baran and Dawber's Diseases of the Nails and their Management}, 5th edition. Wiley-Blackwell, 2019.
\bibitem{nail_discoloration:ref5} Jellinek NJ. "Nail disorders: diagnosis and management." \textit{Medical Clinics of North America}, 2020;104(5):889--904.
\bibitem{nail_discoloration:ref6} Levit EK, Kagen MH, Scher RK. "The ABC rule for clinical detection of subungual melanoma." \textit{Journal of the American Academy of Dermatology}, 2000;42(2):269--274.

\end{thebibliography}
